\begin{center}
\resizebox{0.96\linewidth}{!}{%
\begin{tikzpicture}[
  >=Latex,
  every node/.style={font=\small},
  manifold/.style={very thick},
  tangent/.style={thick,densely dashed,draw=GrisOscuro},
  phasearrow/.style={-{Latex[length=2.2mm]},very thick}
]
  \draw[rounded corners=6pt,thick,GrisOscuro] (0.15,0.15) rectangle (8.25,5.65);
  \node[font=\bfseries] at (4.20,5.30) {Geometría local de una silla hiperbólica};

  \coordinate (e) at (4.05,2.78);
  \draw[tangent] (1.35,0.30) -- (6.75,5.05);
  \draw[tangent] (0.85,4.20) -- (7.40,1.05);

  \draw[manifold,Azul]
    (1.45,0.30) .. controls (2.25,1.10) and (3.35,2.15) .. (e)
    .. controls (4.75,3.48) and (5.70,4.35) .. (6.60,4.90);
  \draw[manifold,Rojo]
    (0.75,4.42) .. controls (2.10,3.78) and (3.18,3.18) .. (e)
    .. controls (4.95,2.38) and (6.05,1.62) .. (7.48,0.85);

  % Sentido del tiempo sobre las ramas estables.
  \draw[phasearrow,Azul] (2.15,0.95) -- (2.88,1.66);
  \draw[phasearrow,Azul] (5.88,4.55) -- (5.13,3.85);
  % Sentido del tiempo sobre las ramas inestables.
  \draw[phasearrow,Rojo] (3.15,3.18) -- (2.28,3.68);
  \draw[phasearrow,Rojo] (4.98,2.35) -- (5.90,1.78);

  \fill[GrisOscuro] (e) circle (2.6pt);
  \node[below right] at (e) {$e$};
  \node[Azul,fill=white,inner sep=2pt] at (2.15,1.90) {$W^s_{\mathrm{loc}}(e)$};
  \node[Rojo,fill=white,inner sep=2pt] at (6.35,2.35) {$W^u_{\mathrm{loc}}(e)$};
  \node[Azul,fill=white,inner sep=2pt] at (6.22,4.52) {$E^s$};
  \node[Rojo,fill=white,inner sep=2pt] at (6.75,1.30) {$E^u$};

  \draw[rounded corners=6pt,thick,GrisOscuro] (8.65,0.15) rectangle (13.95,5.65);
  \node[font=\bfseries] at (11.30,5.30) {Definición dinámica};
  \node[align=left,text width=4.55cm] at (11.30,4.08) {
    \(x\in W^s(e)\) si\\[-1mm]
    \(\varphi_t(x)\to e\) cuando \(t\to+\infty\).};
  \node[align=left,text width=4.55cm] at (11.30,2.78) {
    \(x\in W^u(e)\) si\\[-1mm]
    \(\varphi_t(x)\to e\) cuando \(t\to-\infty\).};
  \draw[GrisOscuro,dashed] (9.10,2.00) -- (13.50,2.00);
  \node[align=center] at (11.30,1.30)
    {$T_eW^s_{\mathrm{loc}}=E^s$,\qquad
     $T_eW^u_{\mathrm{loc}}=E^u$};
  \node[align=center,text=GrisOscuro,font=\scriptsize,text width=4.2cm]
    at (11.30,0.70)
    {tangencia local\\\(\neq\) descripción global};
\end{tikzpicture}%
}

\smallskip
\begin{minipage}{0.92\linewidth}
\small\textbf{Lectura pedagógica.}
Los subespacios espectrales describen las tangentes en el equilibrio; las
variedades invariantes son objetos no lineales. Las flechas azules llegan a la
silla en tiempo futuro y las rojas se alejan de ella. El dibujo es un corte
local bidimensional y no fija la forma global de las separatrices.
\end{minipage}
\end{center}
